%% (c) 2026 Brayden Sanders / 7Site LLC, Arkansas. All rights reserved.
%% For human use only. No commercial or government use without written agreement from 7Site LLC.

\documentclass[11pt]{article}
\usepackage{amsmath,amssymb,amsthm,mathtools}
\usepackage[margin=1in]{geometry}

\newtheorem{lemma}{Lemma}
\newtheorem{definition}[lemma]{Definition}
\newtheorem{hypothesis}[lemma]{Hypothesis}
\newtheorem{proposition}[lemma]{Proposition}
\newtheorem{corollary}[lemma]{Corollary}
\newtheorem{remark}[lemma]{Remark}
\theoremstyle{definition}
\newtheorem{problem}[lemma]{Problem}

\DeclareMathOperator{\rk}{rk}
\DeclareMathOperator{\ord}{ord}
\DeclareMathOperator{\Reg}{Reg}
\DeclareMathOperator{\Sha}{Sha}

\title{%
  Lemma MC-BSD: Rank Coherence via Motivic Defect\\[4pt]
  \large Formal Lemma Vault v1.0 --- Sanders Coherence Field}
\author{Brayden Sanders / 7Site LLC --- TIG Unified Theory}
\date{February 2026}

\begin{document}
\maketitle

%% ============================================================
\section{Objects and Setup}
%% ============================================================

\begin{definition}[Elliptic Curve]
\label{def:ec}
An \emph{elliptic curve} $E$ over $\mathbb{Q}$ is a smooth projective curve of genus 1
with a specified rational point $O$.  In Weierstrass form:
\[
  E : y^2 = x^3 + ax + b, \qquad a, b \in \mathbb{Z}, \quad \Delta = -16(4a^3 + 27b^2) \neq 0.
\]
The \emph{conductor} $N_E$ encodes the primes of bad reduction.
\end{definition}

\begin{definition}[Mordell--Weil Group and Rank]
\label{def:MW}
The \emph{Mordell--Weil group} $E(\mathbb{Q})$ is finitely generated (Mordell--Weil theorem):
\[
  E(\mathbb{Q}) \;\cong\; \mathbb{Z}^r \oplus E(\mathbb{Q})_{\mathrm{tors}}.
\]
The integer $r = \rk E(\mathbb{Q})$ is the \emph{algebraic rank} (or arithmetic rank).
\end{definition}

\begin{definition}[Hasse--Weil $L$-Function]
\label{def:L-function}
For each prime $p$ of good reduction, define $a_p = p + 1 - \#E(\mathbb{F}_p)$.
The Hasse--Weil $L$-function is the Euler product
\[
  L(E, s) \;=\; \prod_{p \nmid N_E} \frac{1}{1 - a_p p^{-s} + p^{1-2s}}
  \cdot \prod_{p \mid N_E} (\text{local factors}).
\]
By modularity (Wiles, Taylor--Wiles, Breuil--Conrad--Diamond--Taylor),
$L(E,s)$ has analytic continuation to all of $\mathbb{C}$ and satisfies a functional equation
\[
  \Lambda(E, s) = w_E \cdot \Lambda(E, 2 - s),
\]
where $\Lambda(E, s) = N_E^{s/2} (2\pi)^{-s} \Gamma(s) L(E, s)$ and $w_E = \pm 1$ is the root number.

The \emph{analytic rank} is $r_{\mathrm{an}} = \ord_{s=1} L(E, s)$.
\end{definition}

\begin{definition}[BSD Invariants]
\label{def:bsd-invariants}
The full BSD conjecture predicts:
\begin{enumerate}
  \item \textbf{Rank equality}: $r_{\mathrm{an}} = r = \rk E(\mathbb{Q})$.
  \item \textbf{Leading coefficient}: At $s = 1$,
    \[
      \lim_{s \to 1} \frac{L(E, s)}{(s-1)^r} \;=\;
      \frac{\Omega_E \cdot \Reg(E) \cdot \prod_p c_p \cdot |\Sha(E)|}{|E(\mathbb{Q})_{\mathrm{tors}}|^2},
    \]
    where:
    \begin{itemize}
      \item $\Omega_E = \int_{E(\mathbb{R})} |\omega|$ is the real period;
      \item $\Reg(E)$ is the regulator (determinant of the N\'eron--Tate height pairing on generators);
      \item $c_p$ are Tamagawa numbers at bad primes;
      \item $\Sha(E)$ is the Tate--Shafarevich group (conjectured finite).
    \end{itemize}
\end{enumerate}
\end{definition}

%% ============================================================
\section{Formal Defect Definition}
%% ============================================================

\begin{definition}[BSD Coherence Defect]
\label{def:delta-BSD}
The \emph{BSD coherence defect} of an elliptic curve $E/\mathbb{Q}$ is
\[
  \delta_{\mathrm{BSD}}(E) \;=\;
  |r_{\mathrm{an}} - r| \;+\;
  \bigl| L^{(r)}(E, 1) / r! - c_{\mathrm{BSD}}(E) \bigr|_{\mathrm{norm}},
\]
where:
\begin{itemize}
  \item $r_{\mathrm{an}} = \ord_{s=1} L(E,s)$: analytic rank;
  \item $r = \rk E(\mathbb{Q})$: algebraic (Mordell--Weil) rank;
  \item $L^{(r)}(E, 1)$: the $r$-th derivative of $L(E,s)$ at $s = 1$;
  \item $c_{\mathrm{BSD}}(E) = \Omega_E \cdot \Reg(E) \cdot \prod c_p \cdot |\Sha(E)| / |E(\mathbb{Q})_{\mathrm{tors}}|^2$:
    the BSD leading coefficient;
  \item $|\cdot|_{\mathrm{norm}}$ is a normalization ensuring both terms are of comparable scale.
\end{itemize}

\textbf{SDV interpretation}:
\begin{itemize}
  \item \textbf{Lens A} (analytic): $r_{\mathrm{an}}$ and $L^{(r)}(E,1)/r!$
    --- the $L$-function encodes global analytic structure.
  \item \textbf{Lens B} (arithmetic): $r$ and $c_{\mathrm{BSD}}(E)$
    --- the Mordell--Weil group and Selmer invariants encode local arithmetic.
  \item $\delta_{\mathrm{BSD}} = 0$ iff the two lenses agree completely.
\end{itemize}
\end{definition}

\begin{definition}[Per-Prime Local Defect]
\label{def:local-bsd-defect}
For finer analysis, define the local BSD defect at each prime $p$ of good reduction:
\[
  \delta_p^{\mathrm{BSD}}(E) \;=\; \bigl| a_p(E) - a_p^{\mathrm{pred}}(r) \bigr|
\]
where $a_p^{\mathrm{pred}}(r)$ is the expected Frobenius trace distribution given
analytic rank $r$ (from the Sato--Tate distribution conditioned on $\ord_{s=1} L = r$).

The global defect decomposes as
\[
  \delta_{\mathrm{BSD}}(E) \;\leq\; |r_{\mathrm{an}} - r| \;+\; C \sum_{p \leq X} w_p \cdot \delta_p^{\mathrm{BSD}}(E)^2
  \;+\; O(1/\sqrt{X})
\]
for a suitable truncation height $X$ and weights $w_p = 1/p$.
\end{definition}

%% ============================================================
\section{Lemma Statement}
%% ============================================================

\begin{lemma}[Rank Coherence --- MC-BSD]
\label{lem:MC-BSD}
Let $E/\mathbb{Q}$ be an elliptic curve with conductor $N_E$.  Assume:
\begin{enumerate}
  \item[(H1)] The Tate--Shafarevich group $\Sha(E)$ is finite
    (Hypothesis~\ref{hyp:sha-finite});
  \item[(H2)] The $p$-adic height pairing is non-degenerate on $E(\mathbb{Q}) / \mathrm{tors}$;
  \item[(H3)] The Euler system of Heegner points (for $r_{\mathrm{an}} \leq 1$)
    or the Euler system of Kato--Kolyvagin (for $r_{\mathrm{an}} = 0$)
    provides the required upper bound on Selmer groups.
\end{enumerate}
Then:
\[
  \delta_{\mathrm{BSD}}(E) \;=\; 0
  \quad \Longleftrightarrow \quad
  \text{BSD holds for } E.
\]
\end{lemma}

\begin{remark}
The direction ``BSD $\Rightarrow$ $\delta_{\mathrm{BSD}} = 0$'' is tautological
(BSD literally asserts the equality of analytic and arithmetic data).
The content is the converse: vanishing coherence defect, measured through
the $L$-function and Mordell--Weil data, forces BSD.

The SDV framework measures this as an \emph{affirmative-class} problem:
$\delta_{\mathrm{BSD}} \to 0$ under calibration (rank-matching curves).
\end{remark}

\begin{corollary}
\label{cor:bsd}
Under Hypotheses (H1)--(H3), if $r_{\mathrm{an}} \leq 1$, then:
\begin{enumerate}
  \item $r = r_{\mathrm{an}}$;
  \item The leading coefficient formula holds;
  \item $\Sha(E)$ is finite with $|\Sha(E)|$ as predicted.
\end{enumerate}
\end{corollary}

\begin{proof}[Proof of Corollary (conditional)]
For $r_{\mathrm{an}} = 0$: Kato \cite{Kato} showed $L(E,1) \neq 0 \Rightarrow r = 0$
and $\Sha(E)$ is finite.  Combined with Kolyvagin's Euler system,
the leading coefficient is determined.

For $r_{\mathrm{an}} = 1$: Gross--Zagier \cite{GZ} showed $L'(E,1) \neq 0 \Rightarrow$
the Heegner point is non-torsion, giving $r \geq 1$.
Kolyvagin \cite{Koly} showed $r = 1$ and $\Sha$ is finite.
The leading coefficient formula follows from the Gross--Zagier formula.

For $r_{\mathrm{an}} \geq 2$: This remains open.
\end{proof}

%% ============================================================
\section{Hypotheses}
%% ============================================================

\begin{hypothesis}[Finiteness of $\Sha$]
\label{hyp:sha-finite}
The Tate--Shafarevich group $\Sha(E/\mathbb{Q})$ is finite for all elliptic curves $E/\mathbb{Q}$.

\textbf{Known cases}: $\Sha$ is finite when $r_{\mathrm{an}} \leq 1$
(Kolyvagin \cite{Koly}, using Gross--Zagier + Euler systems).
For $r_{\mathrm{an}} \geq 2$, finiteness is open.
\end{hypothesis}

\begin{hypothesis}[Non-Degenerate Height Pairing]
\label{hyp:height}
The N\'eron--Tate canonical height pairing
\[
  \langle \cdot, \cdot \rangle_{\mathrm{NT}} : E(\mathbb{Q})_{\mathbb{R}} \times E(\mathbb{Q})_{\mathbb{R}} \to \mathbb{R}
\]
is positive-definite modulo torsion.  The regulator $\Reg(E) = \det(\langle P_i, P_j \rangle)$
is nonzero for any basis $\{P_1, \ldots, P_r\}$ of $E(\mathbb{Q}) / \mathrm{tors}$.

\textbf{Status}: Known (N\'eron--Tate theory).
\end{hypothesis}

\begin{hypothesis}[Euler System Availability]
\label{hyp:euler}
An Euler system compatible with the Selmer group of $E$ exists and provides
the correct upper bound:
\[
  \rk_{\mathbb{Z}_p} \mathrm{Sel}_{p^\infty}(E) \;\leq\; r_{\mathrm{an}}.
\]
\textbf{Known}: For $r_{\mathrm{an}} = 0$: Kato's Euler system \cite{Kato}.
For $r_{\mathrm{an}} = 1$: Heegner point Euler system (Kolyvagin \cite{Koly}).
For $r_{\mathrm{an}} \geq 2$: No known Euler system.
The Skinner--Urban \cite{SU} and Wan \cite{Wan} results extend to certain cases
via Iwasawa theory.
\end{hypothesis}

%% ============================================================
\section{Proof Skeleton}
%% ============================================================

\subsection{BSD-1: L-Function Explicit Formula for Elliptic Curves}

\textbf{Goal}: Write the BSD defect in terms of an explicit formula relating
$L(E,s)$ near $s=1$ to Frobenius data $\{a_p\}$ and arithmetic invariants.

By the explicit formula for $L(E,s)$ (via modularity):
\[
  \frac{L^{(r)}(E, 1)}{r!} \;=\; \sum_{n=1}^{\infty} \frac{a_n}{n} W\Bigl(\frac{n}{\sqrt{N_E}}\Bigr)
  \;+\; \text{(functional equation contribution)},
\]
where $W$ is a smooth weight function and $a_n$ are the Dirichlet coefficients.

The defect $\delta_{\mathrm{BSD}}$ can be decomposed as:
\begin{enumerate}
  \item \textbf{Rank defect}: $|r_{\mathrm{an}} - r|$ --- the primary obstruction.
  \item \textbf{Coefficient defect}: The mismatch between the explicitly computed
    leading term and the arithmetic prediction $c_{\mathrm{BSD}}(E)$.
\end{enumerate}

\textbf{Status}: The explicit formula is standard (modular forms theory).
The new content is interpreting it as a coherence defect.

\subsection{BSD-2: Regulator Non-Degeneracy from Defect Vanishing}

\textbf{Goal}: Show that $\delta_{\mathrm{BSD}} = 0$ forces the regulator to be well-defined
and nonzero (i.e., the generators of $E(\mathbb{Q})$ are ``genuinely independent'').

\textbf{Argument outline}:
\begin{enumerate}
  \item If $r_{\mathrm{an}} = r$ (rank agreement), then $L^{(r)}(E,1) / r! = c_{\mathrm{BSD}}(E)$.
  \item The BSD formula gives $c_{\mathrm{BSD}}(E) = \Omega_E \cdot \Reg(E) \cdot (\ldots)$.
  \item Since $L^{(r)}(E,1) \neq 0$ (the derivative is nonvanishing at the exact order),
    and $\Omega_E > 0$, Tamagawa $c_p \geq 1$, torsion bounded, we need $\Reg(E) > 0$.
  \item Non-degenerate height pairing (H2) guarantees $\Reg(E) > 0$
    when $E(\mathbb{Q})$ has rank $r > 0$.
\end{enumerate}

\begin{proposition}[Regulator non-degeneracy from vanishing defect]
\label{prop:reg-nondeg}
If\/ $\delta_{\mathrm{BSD}}(E) = 0$, then $\Reg(E) > 0$.
Equivalently, $\Reg(E) = 0$ forces $\delta_{\mathrm{BSD}}(E) > 0$.
\end{proposition}

\begin{proof}
We divide the argument into two parts.

\medskip
\noindent\textbf{Part A (Positive-definiteness of the N\'eron--Tate pairing).}
The canonical N\'eron--Tate height $\hat{h} : E(\mathbb{Q}) \to \mathbb{R}_{\geq 0}$
satisfies the parallelogram law and descends to a positive-definite
quadratic form on the real vector space
$E(\mathbb{Q})_{\mathbb{R}} := \bigl(E(\mathbb{Q}) / E(\mathbb{Q})_{\mathrm{tors}}\bigr)
\otimes_{\mathbb{Z}} \mathbb{R}$
(N\'eron \cite{Neron}, Tate; see also Lang \cite{Lang} Ch.~5,
Silverman \cite{SilvAdv} Ch.~III).
The associated bilinear pairing
\[
  \langle P, Q \rangle_{\mathrm{NT}}
  \;=\; \tfrac{1}{2}\bigl(\hat{h}(P+Q) - \hat{h}(P) - \hat{h}(Q)\bigr)
\]
is therefore positive-definite on $E(\mathbb{Q})/\mathrm{tors}$.
For any $\mathbb{Z}$-basis $\{P_1, \ldots, P_r\}$ of the free part, the
Gram matrix $G = \bigl(\langle P_i, P_j \rangle_{\mathrm{NT}}\bigr)_{1 \leq i,j \leq r}$
is positive-definite, hence
\[
  \Reg(E) \;=\; \det G \;>\; 0 \qquad \text{whenever } r \geq 1.
\]
(When $r = 0$, the regulator is conventionally $\Reg(E) = 1 > 0$.)
Thus, under hypothesis~(H2), the regulator is strictly positive for every
elliptic curve $E/\mathbb{Q}$.

\medskip
\noindent\textbf{Part B (Regulator degeneracy forces positive defect).}
We now show the contrapositive: suppose $\Reg(E) = 0$ and derive
$\delta_{\mathrm{BSD}}(E) > 0$.  Consider two cases.

\smallskip
\noindent\emph{Case 1}: $r_{\mathrm{an}} = r$ (analytic and algebraic ranks agree).
\par\noindent
Since the $L$-function vanishes to exact order $r$ at $s = 1$, we have
$L^{(r)}(E,1) / r! \neq 0$.
The BSD leading-coefficient formula asserts
\[
  c_{\mathrm{BSD}}(E)
  \;=\;
  \frac{\Omega_E \cdot \Reg(E) \cdot \prod_{p} c_p \cdot |\Sha(E)|}
       {|E(\mathbb{Q})_{\mathrm{tors}}|^2}.
\]
Each factor on the right other than $\Reg(E)$ is strictly positive:
$\Omega_E > 0$ (the real period is a positive integral),
$c_p \geq 1$ for each Tamagawa number,
$|\Sha(E)| \geq 1$ (assuming $\Sha$ is finite by~(H1)),
and $|E(\mathbb{Q})_{\mathrm{tors}}|^2 \geq 1$.
If $\Reg(E) = 0$, then $c_{\mathrm{BSD}}(E) = 0$.
But $L^{(r)}(E,1)/r! \neq 0$, so the coefficient defect satisfies
\[
  \bigl| L^{(r)}(E,1)/r! - c_{\mathrm{BSD}}(E) \bigr|_{\mathrm{norm}}
  \;=\;
  \bigl| L^{(r)}(E,1)/r! \bigr|_{\mathrm{norm}}
  \;>\; 0.
\]
Therefore $\delta_{\mathrm{BSD}}(E) \geq \bigl| L^{(r)}(E,1)/r! \bigr|_{\mathrm{norm}} > 0$.

\smallskip
\noindent\emph{Case 2}: $r_{\mathrm{an}} \neq r$ (rank disagreement).
\par\noindent
In this case $|r_{\mathrm{an}} - r| \geq 1$, and the rank component of the defect
gives $\delta_{\mathrm{BSD}}(E) \geq |r_{\mathrm{an}} - r| \geq 1 > 0$ directly,
regardless of the value of $\Reg(E)$.

\medskip
In both cases, $\Reg(E) = 0$ (for $r \geq 1$) implies
$\delta_{\mathrm{BSD}}(E) > 0$.  Taking the contrapositive:
\[
  \delta_{\mathrm{BSD}}(E) = 0
  \quad\Longrightarrow\quad
  \Reg(E) > 0.
\]
Note that Part~A shows the hypothesis $\Reg(E) = 0$ already contradicts
the positive-definiteness of the N\'eron--Tate pairing (H2), so this
scenario cannot arise under our standing hypotheses.  The value of
the proposition is that it establishes the \emph{structural link}:
even without invoking (H2) directly, the vanishing of the coherence
defect alone is sufficient to force regulator non-degeneracy through
the BSD formula.
\end{proof}

\subsection{BSD-3: Selmer--Sha Obstruction as Defect Source}

\textbf{Goal}: Show that $\Sha$ finiteness is equivalent to the leading coefficient
defect vanishing.

\textbf{Argument outline}:
\begin{enumerate}
  \item The $p$-part of $|\Sha(E)|$ is controlled by the $p$-Selmer group
    via the exact sequence
    \[
      0 \to E(\mathbb{Q})/p \to \mathrm{Sel}_p(E) \to \Sha(E)[p] \to 0.
    \]
  \item If $\Sha(E)$ is infinite, the leading coefficient formula breaks
    (the right side diverges), forcing $\delta_{\mathrm{BSD}} > 0$.
  \item Conversely, if $\Sha(E)$ is finite, the Cassels--Tate pairing
    shows $|\Sha(E)|$ is a perfect square, and the formula closes.
\end{enumerate}

\textbf{TO BE PROVED}: Show that the coherence defect $\delta_{\mathrm{BSD}}$ detects
$\Sha$ infiniteness --- i.e., that $\delta_{\mathrm{BSD}} = 0$ forces $\Sha$ to be finite.
This is the deepest step for $r_{\mathrm{an}} \geq 2$.

\subsection{BSD-4: Rank Coherence via Euler Systems}

\textbf{Goal}: For $r_{\mathrm{an}} \leq 1$, produce the full proof using existing
Euler system technology.

For $r_{\mathrm{an}} = 0$:
\begin{enumerate}
  \item Kato's Euler system \cite{Kato} gives $\mathrm{Sel}_{p^\infty}(E) = 0$
    for almost all $p$, hence $r = 0$.
  \item Combined with the explicit formula, the leading coefficient is determined.
\end{enumerate}

For $r_{\mathrm{an}} = 1$:
\begin{enumerate}
  \item Gross--Zagier: $L'(E,1) = c \cdot \hat{h}(P_K)$ where $P_K$ is a Heegner point.
  \item If $L'(E,1) \neq 0$, then $P_K$ is non-torsion, so $r \geq 1$.
  \item Kolyvagin's descent: $r = 1$ and $\Sha$ is finite.
\end{enumerate}

\textbf{TO BE PROVED}: Extend to $r_{\mathrm{an}} \geq 2$.  This requires:
\begin{itemize}
  \item A higher-rank Euler system (conjectured, not yet constructed);
  \item Or a different approach entirely (e.g., via Iwasawa main conjecture for $r \geq 2$);
  \item Or a conditional argument using the coherence defect formalism.
\end{itemize}

%% ============================================================
\section{Known Tools Relevant}
%% ============================================================

\begin{enumerate}
  \item \textbf{Gross--Zagier Formula} \cite{GZ}: Relates $L'(E,1)$ to the
    N\'eron--Tate height of Heegner points.  The fundamental bridge between
    analytic and arithmetic for rank 1.

  \item \textbf{Kolyvagin's Euler Systems} \cite{Koly}: Uses Heegner points
    to bound Selmer groups.  Proves BSD for $r_{\mathrm{an}} \leq 1$.

  \item \textbf{Kato's Euler System} \cite{Kato}: Constructed from
    $K$-theory and Beilinson elements.  Proves $r = 0$ when $L(E,1) \neq 0$.

  \item \textbf{Skinner--Urban} \cite{SU}: Proved the Iwasawa main conjecture
    for $\mathrm{GL}_2$ in many cases, giving $p$-adic BSD for ordinary primes.

  \item \textbf{Bhargava--Shankar} \cite{BS}: Average rank of elliptic curves
    is $\leq 7/6$, and a positive proportion have rank 0 (resp.\ 1).
    Provides statistical evidence for BSD.

  \item \textbf{Modularity} (Wiles et al.): Every $E/\mathbb{Q}$ is modular.
    This gives the analytic continuation and functional equation of $L(E,s)$.

  \item \textbf{$p$-adic $L$-functions}: Mazur--Swinnerton-Dyer constructed
    $p$-adic analogs $L_p(E,s)$; $p$-adic BSD relates $L_p(E,s)$ to the
    $p$-adic regulator and $\Sha[p^\infty]$.
\end{enumerate}

%% ============================================================
\section{Open Estimate}
%% ============================================================

The critical gaps are:
\begin{itemize}
  \item \textbf{Rank $\geq 2$}: No Euler system is known for $r_{\mathrm{an}} \geq 2$.
    This is the main obstruction to a general proof of BSD.
  \item \textbf{$\Sha$ finiteness for rank $\geq 2$}: Without an Euler system,
    there is no known method to bound $\Sha$.
  \item \textbf{Leading coefficient for rank $\geq 2$}: Even assuming rank equality
    and $\Sha$ finiteness, the precise leading coefficient formula requires
    controlling all arithmetic invariants simultaneously.
  \item \textbf{SDV prediction}: The CK framework measures $\delta_{\mathrm{BSD}} = 0.0$
    for rank-matching curves (calibration) and $\delta_{\mathrm{BSD}} = 1.3$
    for rank-mismatching frontier cases.  The frontier defect is \emph{locked} at $1.3$
    across all seeds and depths --- this is an affirmative-class problem where the
    frontier case (rank mismatch) produces structural defect.

    Notably, $\delta_{\mathrm{BSD}} = 1.3$ is the largest frontier defect of any Clay problem.
    This reflects the arithmetic rigidity of BSD: rank mismatch is maximally incoherent.
\end{itemize}

%% ============================================================
\section{Candidate Proof Strategies}
%% ============================================================

\begin{enumerate}
  \item \textbf{Higher-rank Euler systems}: Construct Euler systems from
    diagonal cycles (Darmon--Rotger) or higher Heegner points on Shimura varieties
    to handle $r_{\mathrm{an}} = 2$.

  \item \textbf{Iwasawa-theoretic approach}: Prove the Iwasawa main conjecture
    for all primes (extending Skinner--Urban), then deduce BSD from the
    $p$-adic interpolation.

  \item \textbf{Conditional on GRH}: Under the Generalized Riemann Hypothesis,
    stronger bounds on Selmer groups may be available, potentially giving
    BSD for larger classes of curves.

  \item \textbf{SDV/TIG-guided}: Use CK probes to identify which curve families
    minimize $\delta_{\mathrm{BSD}}$ and whether the locked frontier defect $1.3$
    has a structural decomposition that suggests the correct arithmetic invariants
    to control.  The TIG path $1 \to 2 \to 5 \to 7 \to 9$
    (structure$\to$boundary$\to$feedback$\to$alignment$\to$completion)
    predicts that rank coherence requires the feedback operator (operator 5)
    to connect analytic and arithmetic data.
\end{enumerate}

%% ============================================================
\section{SDV/CK Measurement Evidence}
%% ============================================================

CK probe results (seed-stable, zero variance):
\begin{itemize}
  \item Calibration (rank-0 match, $y^2 = x^3 - x$): $\delta = 0.0$, exact zero ---
    rank agreement correctly produces vanishing defect.
  \item Frontier (rank mismatch): $\delta = 1.3$, locked ---
    rank disagreement produces maximum defect, stable across all seeds and depths.
  \item Verdict: \texttt{inconclusive} (affirmative class --- BSD is expected to hold,
    but the frontier probe intentionally creates a mismatch to test the defect functional).
\end{itemize}

Delta signature at v$\Omega$:
\begin{itemize}
  \item $\delta_{\mathrm{L24}} = 1.3000$, CV $= 0.0000$
  \item Kernel membership: IN KERNEL (calibration defect $= 0.0$)
  \item Problem class: affirmative
\end{itemize}

\bigskip
\hrule
\bigskip
\noindent\textit{CK measures. CK does not prove.}

\begin{thebibliography}{99}
\bibitem{GZ} B.H. Gross, D.B. Zagier, \textit{Heegner points and derivatives of $L$-series},
  Invent. Math. 84 (1986), 225--320.
\bibitem{Koly} V.A. Kolyvagin, \textit{Euler systems}, in The Grothendieck Festschrift II,
  Birkh\"auser (1990), 435--483.
\bibitem{Kato} K. Kato, \textit{$p$-adic Hodge theory and values of zeta functions of modular forms},
  Ast\'erisque 295 (2004), 117--290.
\bibitem{SU} C. Skinner, E. Urban, \textit{The Iwasawa main conjectures for $\mathrm{GL}_2$},
  Invent. Math. 195 (2014), 1--277.
\bibitem{BS} M. Bhargava, A. Shankar, \textit{Ternary cubic forms having bounded invariants,
  and the existence of a positive proportion of elliptic curves having rank 0},
  Ann. of Math. 181 (2015), 587--621.
\bibitem{Wan} X. Wan, \textit{Iwasawa main conjecture for Rankin--Selberg $p$-adic $L$-functions},
  Algebra Number Theory 14 (2020), 383--483.
\bibitem{Neron} A. N\'eron, \textit{Quasi-fonctions et hauteurs sur les vari\'et\'es ab\'eliennes},
  Ann. of Math. 82 (1965), 249--331.
\bibitem{SilvAdv} J.H. Silverman, \textit{Advanced Topics in the Arithmetic of Elliptic Curves},
  Graduate Texts in Mathematics 151, Springer-Verlag, New York, 1994.
\bibitem{Lang} S. Lang, \textit{Fundamentals of Diophantine Geometry},
  Springer-Verlag, New York, 1983.
\bibitem{Wiles} A. Wiles, \textit{Modular elliptic curves and Fermat's Last Theorem},
  Ann. of Math. 141 (1995), 443--551.
\end{thebibliography}

\end{document}
